\subsection{Native validation and measured limits}

The comparison pins the preceding maintained producer and persistent checker
at commit \texttt{4360509f450e} (the full hash is in each data file). Only the
existing producer selector and compressed-search dispatcher differ among
preexisting runtime modules; the new private producer is additive. Every
preexisting verifier, including the complete compressed and literal source
checkers, is byte-identical. The driver freezes 434 Python/corpus source
hashes before each run and checks them again afterward.

All 1,116 maintained tests pass with Regina available, in 222.967 seconds.
The four new methods compare 350 random greedy presentations with ordinary
updates and literal batch replay; exercise a sixteen-generator binary-power
chain with $2^{128}$ repetitions; test summary-cache invalidation, cancellation,
work/node exhaustion and atomic publication; and check both the terminal
fast path and an actual two-batch source certificate. The binary-capacity test
disables expansion, equality, reduction and inverse helpers during export
and independent replay. It is a supplied transformation schedule, not evidence
that greedy discovery chooses that schedule on a knot diagram.

The separate pinned audit checks 84 diagrams, including the maintained corpus,
random braid closures and stabilized circles. All default and optional-batch
certificates match the predecessor, with no gained or lost positives. There
are 840 original-source literal/compressed replays across the two packages,
and 24 searches activate persistent production. Activation includes searches
that later stall; it does not imply a newly accepted proof. The driver also
retains the 400 abstract DAG cases, 1,600 compressed replays and 800 literal
replays used to check the unchanged batch interface. The audit takes 8.040
seconds; its source hashes and complete evidence are retained separately
from timings.

\paragraph{Measurement protocol.}
All timings run serially after tests and audit completion. Four arms comprise
the predecessor, the current producer and an identical A/A control of each.
Each case has one excluded warmup and five measured rounds in shuffled order.
Every timed call constructs and validates a fresh source diagram and includes
the host's mandatory independent replay. The three experiments total 520
measured calls and 104 warmups, all complete. Ratios below are medians of
within-round old/new ratios; values above one favor the current producer.
They are not ratios of displayed median times. These small, single-environment
measurements do not establish a broad recognition speedup.

\paragraph{An activating source-bound group call.}
The actual \texttt{survivor-10} source discovers two batches of seven and two
generators, then normalizes and reaches the existing terminal. Its unchanged
726-byte certificate is replayed inside the timed group call and additionally
checked by the literal verifier outside timing. One private producer block
uses 57 nodes, 22 context siblings, maximum concatenation depth three and a
combined private/ordinary peak of 158 nodes. Final ordinary producer nodes
fall from 177 to 127, while total charged search work, including source
recovery, is essentially unchanged: 5,464 versus 5,467. The timing likewise
shows no material gain:
\begin{center}\small
\begin{tabular}{lrrrrr}
\toprule
Case & Old ms & New ms & Old/new & Old A/A & New A/A \\
\midrule
survivor-10 & 3.747 & 3.686 & 1.017 & 1.011 & 0.992 \\
\bottomrule
\end{tabular}
\end{center}


The private-node peak concerns the raw block; final ordinary-node counts also
include subsequent terminal/normalization work. They are not interchangeable
memory measurements, and neither includes Python container and integer bytes.

\paragraph{Terminal fast-path controls.}
Stabilized-circle group calls finish in one batch and therefore avoid the
private import. Their paired ratios range from 0.981 to 1.012; the current
A/A controls range from 0.993 to 1.014. This is an overhead control, not a
performance demonstration of persistent production:
\begin{center}\small
\begin{tabular}{lrrrrr}
\toprule
Case & Old ms & New ms & Old/new & Old A/A & New A/A \\
\midrule
circle-8 & 1.253 & 1.253 & 1.000 & 0.992 & 0.998 \\
circle-16 & 2.515 & 2.510 & 1.012 & 0.993 & 1.008 \\
circle-32 & 4.743 & 4.902 & 0.981 & 0.991 & 1.000 \\
circle-64 & 9.555 & 9.578 & 1.005 & 0.988 & 0.998 \\
circle-128 & 19.991 & 20.140 & 0.995 & 0.994 & 1.014 \\
circle-256 & 38.915 & 39.168 & 0.991 & 0.983 & 0.993 \\
\bottomrule
\end{tabular}
\end{center}


\paragraph{Complete native recognition.}
The nineteen-diagram corpus runs the same optional batch portfolio in both
engines, including the bounded speculative trial and complete fallback.
All 380 measured calls and 76 warmups complete. Ratios span 0.944--1.026,
with several small regressions. Old A/A ratios span 0.957--1.020 and new A/A
ratios span 0.982--1.015. There is no broad ordinary-workload gain:
\begin{center}\small
\begin{tabular}{lrrrrr}
\toprule
Case & Old ms & New ms & Old/new & Old A/A & New A/A \\
\midrule
survivor-00 & 13.957 & 14.511 & 0.958 & 0.996 & 0.999 \\
survivor-01 & 28.418 & 28.902 & 0.973 & 0.957 & 0.982 \\
survivor-02 & 9.260 & 9.563 & 0.965 & 1.002 & 1.003 \\
survivor-03 & 38.839 & 39.286 & 0.993 & 1.000 & 1.011 \\
mirror-03 & 5.741 & 5.758 & 0.995 & 0.997 & 1.010 \\
survivor-04 & 29.682 & 30.025 & 0.992 & 0.993 & 0.999 \\
survivor-05 & 15.225 & 15.677 & 0.969 & 0.978 & 1.000 \\
survivor-06 & 5.098 & 5.172 & 0.986 & 0.998 & 0.997 \\
survivor-07 & 20.408 & 19.805 & 1.026 & 1.018 & 0.995 \\
survivor-08 & 6.867 & 7.166 & 0.958 & 1.008 & 1.003 \\
mirror-08 & 15.169 & 15.522 & 0.979 & 1.005 & 1.001 \\
survivor-09 & 5.493 & 5.823 & 0.944 & 1.017 & 1.015 \\
survivor-10 & 4.937 & 5.023 & 0.984 & 1.002 & 1.006 \\
survivor-11 & 4.665 & 4.663 & 1.000 & 1.002 & 0.996 \\
gordian & 1261.500 & 1253.230 & 1.011 & 1.020 & 0.987 \\
trefoil & 0.054 & 0.055 & 0.990 & 0.996 & 0.996 \\
figure\_eight & 0.061 & 0.061 & 0.998 & 1.008 & 1.007 \\
conway & 0.971 & 0.973 & 0.997 & 0.998 & 0.996 \\
kinoshita\_terasaka & 1.002 & 1.002 & 0.996 & 0.986 & 1.002 \\
\bottomrule
\end{tabular}
\end{center}


For \texttt{survivor-10}, the complete pipeline ratio is 0.984, despite the
smaller producer grammar. Gordian's trial still exhausts 89,905 charged units,
comprising 39,904 recovery units and the 50,000 post-recovery allowance plus
the failing charge. Both trials perform two batches eliminating 127 generators
and then fall back. Persistent production reduces ordinary trial nodes from
2,788 to 2,122, with 735 private nodes and a raw-block combined peak of 1,982.
The final accepted fallback certificate is unchanged; the paired whole-call
ratio 1.011 is not evidence that the trial now solves Gordian.

The implementation is \path{../fast/fastunknot/persistent_elimination.py};
reproduction is \path{../fast/compressed_word_research/persistent_producer.py}.
Complete JSON, logs and table generation are retained as
\path{data/persistent-producer-*} and \path{data/persistent_producer_tables.py}.
The earlier checker kernel ratios concern supplied replay schedules and are
not reused as producer speedups. The established improvement here is the
whole-block discovery bound and reduced intermediate grammar allocation on
the activating sources, with explicit ordinary-workload overhead retained.
